\section{Spectral and arithmetic laws of exact response memory}\label{sec:r69-spectral}
The instrument realization theorem preserves every terminal effect, whereas a prescribed task can require substantially less memory. We now determine the exact memory problem for a whole family of single-frequency protected responses, not just one damping factor or the previously solved three-label case. The criterion reduces the problem to classical stochastic eigenvalue regions. A self-contained sector argument and elementary rational approximation then give an exact worst-phase law and matching growth at every fixed badly approximable phase. None of the spectral-region, positive-realization, or sector machinery is claimed as newly invented \cite{DmitrievDynkin46,JohnsonPaparella17,KirklandSmigoc22}. The contributions considered here are their exact response-level equivalence with a prescribed initial label, their consequences for the physical squared task, and the phase-sensitive matched resource law.

\subsection{A physical family and the exact spectral criterion}
Let $0<\alpha\le1$, $0<\gamma\le1$, and $\vartheta,\varphi\in\mathbb R$, with $\sin\vartheta\ne0$. Consider the binary response
\begin{equation}\label{eq:r69-response}
 b_n=\frac12\bigl[1+\gamma\alpha^n\cos(n\vartheta+\varphi)\bigr],
 \qquad n\ge0.
\end{equation}
The rotation, contrast and preparation phase specify one prepared instance; independence of the minimum size does not assert one fixed decoder for unknown phases or contrasts. It is physically prepared on a qubit: take the Bloch vector
$\gamma(\sin\varphi,0,\cos\varphi)$, the channel
\[
 \Phi(X)=\alpha U_\vartheta XU_\vartheta^*
               +(1-\alpha)\tr(X)I/2,\qquad
 U_\vartheta=\begin{pmatrix}\cos(\vartheta/2)&-\sin(\vartheta/2)\\
                  \sin(\vartheta/2)&\cos(\vartheta/2)\end{pmatrix},
\]
and terminal effect $E=(I+\sigma_z)/2$. Its Born probabilities are exactly \eqref{eq:r69-response}. Theorem~\ref{thm:r67-embedding} embeds this response into the same three-outcome squared-loss experiment, with protected dimension four, fixed decoder dictionary $\{-1,1\}$, and least expected cost $1/p$. It forces the probability of output $+1$ after $n$ failures to equal $b_n$. Thus the mathematical realization problem below is the exact optimal-observer problem of an actual bounded squared task, not a sequence supplied by a loss oracle.

Write $m(\alpha,\vartheta;\gamma,\varphi)$ for the least $M$ for which
\[
 b_n=e_1^{\mathsf T}Q^nd\quad(n\ge0),\qquad
 Q\ge0,\quad Q\mathbf1=\mathbf1,\quad d\in[0,1]^M.
\]
The initial label is prescribed and counted. Let
\[
 \mathcal K_M=\{z\in\mathbb C:z\text{ is an eigenvalue of some }
                         M\times M\text{ row-stochastic matrix}\}.
\]
These are the classical Karpelevi\v c regions; their full boundary classification is prior work \cite{JohnsonPaparella17}. We will not use that classification as an unproved step in any bound below.

\begin{theorem}[Exact spectral realization, with no extra initial label]\label{thm:r69-spectral}
For the nondegenerate family \eqref{eq:r69-response},
\begin{equation}\label{eq:r69-minimal}
 m(\alpha,\vartheta;\gamma,\varphi)
       =\min\{M:\alpha e^{i\vartheta}\in\mathcal K_M\}.
\end{equation}
The value is independent of every nonzero contrast $\gamma$ and of the preparation phase $\varphi$; denote it by $m(\alpha,\vartheta)$. The criterion is necessary and sufficient over all stochastic matrices, including reducible matrices and arbitrary randomized decoders. For $0<\alpha<1$ it is finite. For $\alpha=1$ and $\vartheta/(2\pi)=a/q$ in lowest terms it is exactly $q$; for irrational $\vartheta/(2\pi)$ it is infinite.

In the squared-task embedding, these are exact live-label counts at cost $1/p$. At each fixed smaller memory bound, the attained risk is strictly above the full-history value $2$. Event-stopped and state-stopped total-label upper counts remain respectively $\max\{M,2\}$ and $M+2$.
\end{theorem}
\begin{proof}
Put $z=\alpha e^{i\vartheta}$. The response has the three distinct exponential modes $1,z,\bar z$, with nonzero coefficients. Cayley--Hamilton for any realizing $Q$ makes its characteristic polynomial an annihilating shift polynomial of $b_n$. Substitution and Vandermonde independence force that polynomial to vanish at $z$. This proves necessity without assuming a minimal linear representation has positive coordinates.

Conversely let $Qv=zv$ with $v\ne0$. Normalize a largest-modulus coordinate to one by multiplying $v$ by one complex scalar, and permute that coordinate into position one. Then $v_1=1$ and $|v_j|\le1$. The column
\[
 d_j=\tfrac12[1+\gamma\operatorname{Re}(e^{i\varphi}v_j)]
\]
lies in $[0,1]^M$. Row normalization and $Q^nv=z^nv$ give $e_1^{\mathsf T}Q^nd=b_n$ for every $n$. The same $M$ labels suffice: there is no preliminary random initialization or extra anchor. Unreachable labels can be pruned. Finiteness for $\alpha<1$ follows from the regular-polygon construction proved below, or from the inherited polygon upper in Theorem~\ref{thm:qubit-positive-phase}.

At $|z|=1$, equality in the triangle inequality at a largest-modulus eigenvector coordinate forces each positive successor to have the same modulus and value $z$ times the current value. Following positive successors reaches a directed cycle of length $l\le M$, so $z^l=1$. An irrational phase is therefore impossible. If $z$ has order $q$, every such cycle has length divisible by $q$, giving $M\ge q$; a $q$-cycle realizes it. Finally the squared embedding forces the whole response. Compact attainment at fixed $M$ was proved in Proposition~\ref{prop:r67-density}; equality of its risk with $2$ would give the excluded exact response. The existing packing lemma gives the terminal counts.
\end{proof}

The nondegeneracy assumptions cannot be suppressed. Zero contrast gives one constant label. At $\alpha=0$, $b_0=(1+\gamma\cos\varphi)/2$ and $b_n=1/2$ for $n\ge1$; one label suffices if $\gamma\cos\varphi=0$, and otherwise two are necessary and sufficient. Real phases and constant cases are classified in the supplement. In particular a vanishing oscillation amplitude is not the limit of an exact-memory lower valid for nonzero amplitude.

For algebraic $x=\operatorname{Re}z$, $y=\operatorname{Im}z\ne0$, spectral membership has the compact polynomial certificate
\begin{equation}\label{eq:r69-bilinear}
 \begin{gathered}
 Q\ge0,\ Q\mathbf1=\mathbf1,\quad
 Qa=xa-yc,\quad Qc=ya+xc,\\
 a_1=1,\ c_1=0,\quad a_j^2+c_j^2\le1\quad(1\le j\le M).
 \end{gathered}
\end{equation}
It has $M^2+2M$ real variables and polynomial degree at most two. Normalization of an eigenvector proves equivalence in both directions. Real-algebraic elimination decides it at each fixed $M$, and the finite polygon upper terminates a minimal-size search for $|z|<1$. This improves the form of the certificate, not the general complexity of real-algebraic elimination. A chosen algebraic preparation and contrast then yield algebraic decoder entries. State labels, encodings of those entries, exact-sampling resources and runtime remain separate accounts.

\subsection{The sector bound, its equality case, and exact finite phases}
We include the classical stochastic sector argument to expose every load-bearing inequality. The equality structure is consistent with the Type-0 classification of Kirkland and \v Smigoc \cite{KirklandSmigoc22}; it is not a new classification of stochastic spectral arcs.

\begin{lemma}[Stochastic sector inequality and rigidity]\label{lem:r69-sector}
If $z=x+iy$, $y\ne0$, is an eigenvalue of an $M\times M$ stochastic matrix, then $M\ge3$ and
\begin{equation}\label{eq:r69-sector}
 |y|\le(1-x)\cot(\pi/M).
\end{equation}
If equality holds, then, up to simultaneous relabelling and reversing the orientation when needed, that matrix is
\begin{equation}\label{eq:r69-lazy-cycle}
 Q=(1-t)I+tC_M,\qquad 0<t\le1,
\end{equation}
where $C_M$ is the directed $M$-cycle. Its eigenvalue on the equality ray is $1-t+te^{2\pi i/M}$, or its conjugate.
\end{lemma}
\begin{proof}
Conjugate the eigenvector when necessary so $y>0$. Let $P$ be the convex hull of its complex coordinates. The matrix
$e^{s(Q-I)}=e^{-s}\sum_{n\ge0}s^nQ^n/n!$ is stochastic for $s\ge0$, so multiplication by $e^{s(z-1)}$ maps $P$ into itself. Since $z$ is nonreal, $1-x>0$, and the spiral of any nonzero coordinate visits every angle, a finite set of points of this spiral surrounds zero. Hence zero lies in the interior of $P$.

List its $l\le M$ vertices counterclockwise as $r_je^{i\theta_j}$, with $r_j>0$. Their consecutive angular gaps $\Delta_j$ belong to $(0,\pi)$ and sum to $2\pi$. Put $u=1-x>0$ and $v=y>0$. The initial velocity $(-u+iv)r_je^{i\theta_j}$ must point into the half-plane to the left of the outgoing edge. Taking its signed two-dimensional cross product with that edge gives
\[
 r_j/r_{j+1}\le \cos\Delta_j+(u/v)\sin\Delta_j.
\]
All right sides are positive. Their product is at least one. Write $\phi=\arctan(u/v)\in(0,\pi/2)$. The function
$\log(\cos\Delta+(u/v)\sin\Delta)
 =\log\sqrt{1+(u/v)^2}+\log\cos(\Delta-\phi)$
is strictly concave on the interval of positive factors in $(0,\pi)$. Jensen's inequality therefore gives
\[
 1\le\prod_j[\cos\Delta_j+(u/v)\sin\Delta_j]
 \le[\cos(2\pi/l)+(u/v)\sin(2\pi/l)]^l.
\]
As $l\ge3$, this implies $u/v\ge\tan(\pi/l)\ge\tan(\pi/M)$, proving \eqref{eq:r69-sector}.

At equality, $l=M$, every gap is $2\pi/M$, and all radius ratios equal one: the polygon is regular and every coordinate is one distinct vertex. The velocity at a vertex is along its outgoing edge. But
$(z-1)v_j=\sum_kQ_{jk}(v_k-v_j)$ is a nonnegative combination of chords. Strict convexity of this polygon makes all chords other than the self-chord and outgoing edge point strictly to the interior side of that edge. Their coefficients must therefore vanish. The remaining coefficient $t=Q_{j,j+1}$ satisfies $z-1=t(e^{2\pi i/M}-1)$ and is independent of $j$. Row normalization gives \eqref{eq:r69-lazy-cycle} and $t\le1$.
\end{proof}

\begin{theorem}[Exact first-sector memory threshold and optimizer]\label{thm:r69-sector-phase}
For $M\ge3$ and $0<\vartheta\le2\pi/M$, a response realization with at most $M$ live labels exists exactly when
\begin{equation}\label{eq:r69-sector-threshold}
 0<\alpha\le r_M(\vartheta)
 :=\frac1{\cos\vartheta+\sin\vartheta\tan(\pi/M)}.
\end{equation}
The threshold is attained. At $\alpha=r_M(\vartheta)$ the least memory is exactly $M$, the transition matrix is the lazy directed cycle in \eqref{eq:r69-lazy-cycle}, and, in cyclic coordinates with prescribed initial label zero, its decoder is uniquely
\[
 d_j=\tfrac12[1+\gamma\cos(\varphi+2\pi j/M)].
\]
This is global over arbitrary stochastic transition and decoder rows, not an optimization restricted in advance to circulants.
\end{theorem}
\begin{proof}
The spectral criterion and \eqref{eq:r69-sector} prove necessity. For sufficiency put $\omega=e^{2\pi i/M}$ and $J=\mathbf1\mathbf1^{\mathsf T}/M$. The nonnegative coefficients
\[
 b=\frac{\alpha\sin\vartheta}{\sin(2\pi/M)},\qquad
 a=\alpha\cos\vartheta-b\cos(2\pi/M),\qquad
 c=1-a-b
\]
sum to one under \eqref{eq:r69-sector-threshold}. The inequalities $a,b\ge0$ follow from $0<\vartheta\le2\pi/M$; $c\ge0$ is exactly the stated threshold. Then $Q=aI+bC_M+cJ$ has eigenvector $(1,\omega,\ldots,\omega^{M-1})$ at eigenvalue $\alpha e^{i\vartheta}$ and the displayed decoder realizes the response.

At the threshold, strict monotonicity of $\cot(\pi/M)$ excludes every smaller size. Lemma~\ref{lem:r69-sector} fixes $Q$. Here $b=t>0$, so $C_M=(Q-(1-t)I)/t$. The rows $e_0^{\mathsf T}Q^n$, $0\le n<M$, span all coordinate rows, because the rows $e_0^{\mathsf T}C_M^j$ do and each is a polynomial combination of the former. Thus all response coefficients determine the decoder uniquely. This proves the rigidity assertion including the endpoint $t=1$.
\end{proof}

For the inherited algebraic angle $\cos\vartheta=3/5$, $\sin\vartheta=4/5$, let
\[
 a_j=\frac5{3+4\tan(\pi/j)},\qquad j=3,4,5,6.
\]
One has $0<\vartheta<2\pi/6$. Consequently exact memory is three on $(0,a_3]$, four on $(a_3,a_4]$, five on $(a_4,a_5]$, and six on $(a_5,a_6]$. In particular $a_4=5/7$. Above $a_6$ at least seven labels are needed; no assertion that a regular polygon is always globally optimal beyond the first sector is made. The old three-label threshold is recovered rather than reproved as a new independent result.

\subsection{Exact worst-phase size and arithmetic improvements}
For an integer $q\ge3$ define the distance to the nearest $q$th root angle by
\[
 \delta_q(\vartheta)=\min_{j\in\mathbb Z}|\vartheta-2\pi j/q|
             \quad\text{modulo }2\pi,\qquad 0\le\delta_q\le\pi/q,
\]
and define the radial boundary of the regular polygon by
\begin{equation}\label{eq:r69-polygon-radius}
 R_q(\vartheta)=\frac{\cos(\pi/q)}{\cos(\pi/q-\delta_q(\vartheta))}.
\end{equation}

\begin{lemma}[Directional polygon construction and power obstruction]\label{lem:r69-arithmetic-tools}
If $0<\alpha\le R_q(\vartheta)$, a $q$-label exact response realization exists. Conversely, any $M$-label realization satisfies, for every positive integer $k$,
\begin{equation}\label{eq:r69-power-obstruction}
 \alpha^k\bigl[\cos(k\vartheta)+|\sin(k\vartheta)|\tan(\pi/M)\bigr]\le1.
\end{equation}
For $h=\pi/q$ and $0\le\delta=\delta_q\le h$,
\begin{equation}\label{eq:r69-polygon-error}
 1-R_q(\vartheta)\le2h\delta.
\end{equation}
\end{lemma}
\begin{proof}
The intersection of the ray at angle $\vartheta$ with the chord joining its two adjacent unit roots has radius \eqref{eq:r69-polygon-radius}. Thus $\alpha e^{i\vartheta}$ is in their triangle with zero. Write it as $a\omega^j+b\omega^{j+1}+c0$, with $a,b,c\ge0$ and sum one. The stochastic matrix $aC_q^j+bC_q^{j+1}+cJ$ has that eigenvalue, giving the first assertion by Theorem~\ref{thm:r69-spectral}. The matrix $Q^k$ of any realization is still stochastic and has eigenvalue $\alpha^ke^{ik\vartheta}$. Apply the sector lemma; for a real eigenvalue the displayed inequality is automatic. Finally $\cos(h-\delta)\ge\cos h\ge1/2$ and
$\cos(h-\delta)-\cos h=2\sin(h-\delta/2)\sin(\delta/2)\le h\delta$.
Division proves \eqref{eq:r69-polygon-error}.
\end{proof}

\begin{theorem}[Exact worst-phase memory]\label{thm:r69-worst-phase}
For every $0<\alpha<1$,
\begin{equation}\label{eq:r69-worst-phase}
 \max_{\sin\vartheta\ne0}m(\alpha,\vartheta)
  =\max\left\{3,\left\lceil\frac\pi{\arccos\alpha}\right\rceil\right\}.
\end{equation}
In particular the worst-phase exact label count is asymptotic to
$\pi/\sqrt{2(1-\alpha)}$ as $\alpha\uparrow1$. The phase may depend on $\alpha$; this is not a lower rate for each fixed phase.
\end{theorem}
\begin{proof}
Let the displayed integer be $K$. The regular $K$-gon contains the disk of radius $\cos(\pi/K)\ge\alpha$, so the polygon lemma gives the upper simultaneously for all phases. If $K=3$, three distinct nonzero modes force at least three states. Otherwise put $M=K-1\ge3$ and choose $\vartheta=\pi/M$. The definition of $K$ gives $\alpha>\cos(\pi/M)$; the exact first-sector threshold at this angle is $\cos(\pi/M)$. No $M$-label realization exists, proving the lower and attainment of the maximum. The asymptotic follows from $1-\cos t\sim t^2/2$.
\end{proof}

A number $x$ is called badly approximable here with constant $c>0$ when
\begin{equation}\label{eq:r69-diophantine}
 \|kx\|_{\mathbb R/\mathbb Z}\ge c/k\qquad(k\ge1).
\end{equation}
The constant is part of the hypothesis, not something inferred from a finite list of continued-fraction coefficients. The next theorem holds for every such phase and proves a different, matching fixed-phase rate.

\begin{theorem}[Sharp cubic-root memory at a fixed Diophantine phase]\label{thm:r69-cubic}
Let $x=\vartheta/(2\pi)$ satisfy \eqref{eq:r69-diophantine}. Then, for $0<\alpha<1$ and $\epsilon=1-\alpha$,
\begin{equation}\label{eq:r69-cubic-bounds}
 \left(\frac{\pi^2c}{16\epsilon}\right)^{1/3}
 \le m(\alpha,\vartheta)
 \le \max\left\{4,\left\lceil\frac3c\right\rceil,
              \left\lceil\left(\frac{4\pi^2}{c^2\epsilon}\right)^{1/3}\right\rceil\right\}.
\end{equation}
Thus $m(\alpha,\vartheta)=\Theta_c((1-\alpha)^{-1/3})$. The exact live-state bits are $\frac13\log_2(1/(1-\alpha))+O_c(1)$ in the ideal finite-label model. These are exact stochastic realizations at the original least expected-cost budget, not approximation of the response or additional physical observations. The statement does not supply matching program-bit or exact-sampling complexity bounds.
\end{theorem}
\begin{proof}
We use only the following elementary approximation fact. For any irrational $x$ and integer $Q\ge1$, among the $Q+1$ fractional parts $0,x,\ldots,Qx$ two fall in the same one of $Q$ equal intervals. Subtracting them gives $1\le k\le Q$ with $0<\|kx\|\le1/Q$.

For the lower, take a realizing size $M\ge3$ and apply that fact with $Q=2M$. Set $\delta=2\pi\|kx\|$ and $h=\pi/M$. Then
\[
 \pi c/M\le\delta\le h\le\pi/3.
\]
Writing $F=\cos\delta+\sin\delta\tan h$, we have $F\le\sec h\le2$. Also $\sin\delta\ge3\delta/4$ on $[0,\pi/3]$, $\tan h\ge h$ and $1-\cos\delta\le\delta^2/2$, so
\[
 F-1\ge3h\delta/4-\delta^2/2\ge h\delta/4.
\]
By \eqref{eq:r69-power-obstruction}, $\alpha^kF\le1$. Hence
\[
 2M\epsilon\ge k\epsilon\ge1-\alpha^k
       \ge1-1/F\ge h\delta/8
       \ge\pi^2c/(8M^2).
\]
This proves the lower constant.

For the upper, let $Q$ be the integer on the right of \eqref{eq:r69-cubic-bounds}. The approximation fact provides integers $1\le q\le Q$ and $a$ with $|qx-a|\le1/Q$. The lower bound \eqref{eq:r69-diophantine} forces $q\ge cQ\ge3$. The distance from $\vartheta$ to a $q$th-root angle is at most $2\pi/(qQ)$. The polygon-error bound gives
\[
 1-R_q(\vartheta)\le\frac{4\pi^2}{q^2Q}
       \le\frac{4\pi^2}{c^2Q^3}\le\epsilon.
\]
Thus $\alpha\le R_q(\vartheta)$ and the explicit circulant construction realizes the response with $q\le Q$ labels. Both bounds apply for the same fixed phase as $\epsilon\downarrow0$. Taking base-two logarithms and rounding the label bits proves the last assertion.
\end{proof}

The hypothesis is nonempty and has elementary physical examples. For
$x=(3-\sqrt5)/2$, choose the closest integer $a$ to $kx$. The integer
$a^2-3ak+k^2$ is nonzero. Since the other root is $x'=(3+\sqrt5)/2$ and
$|a/k-x'|\le\sqrt5+1/2<3$, factorization gives
$|a/k-x|>1/(3k^2)$. Thus \eqref{eq:r69-diophantine} holds with $c=1/3$. This provides a fixed qubit rotation with the matched exponent. Its trigonometric primitive entries are not asserted to be algebraic. In particular the Diophantine condition has not been proved here for the different algebraic rotation $(3+4i)/5$.

\begin{corollary}[A stronger fixed-algebraic-angle exclusion]\label{cor:r69-algebraic-angle}
For the inherited angle $e^{i\vartheta}=(3+4i)/5$ and $0<\alpha<1$, every exact $M$-label stochastic realization satisfies
\begin{equation}\label{eq:r69-algebraic-bound}
 1-\alpha\ge\frac3{16M^2\,25^M}.
\end{equation}
Consequently its exact memory is at least of order $\log(1/(1-\alpha))$ as $\alpha\uparrow1$. This improves the earlier $\log/\log\log$ lower, without assigning the cubic-root law to this angle or claiming a matching upper for it.
\end{corollary}
\begin{proof}
Again choose $1\le k\le2M$ with $\delta=2\pi\|k\vartheta/(2\pi)\|\le\pi/M$. The Gaussian integer $(3+4i)^k-5^k$ is nonzero: otherwise the unit complex number would be a root of unity, forcing $6/5$ to be a rational algebraic integer. Its modulus is at least one. Therefore
$\delta\ge|1-e^{ik\vartheta}|\ge5^{-k}\ge25^{-M}$.
Repeat the lower proof of Theorem~\ref{thm:r69-cubic} up to $2M(1-\alpha)\ge h\delta/8$, with $h=\pi/M$. This gives the stronger constant $\pi/(16M^2 25^M)$ and hence \eqref{eq:r69-algebraic-bound}. Taking logarithms gives
$\log(1/(1-\alpha))\le\log(16/3)+2\log M+M\log25$, proving the order assertion.
\end{proof}

\subsection{Finite spectral mixtures and higher-dimensional protected systems}
The matching exponent is not confined to a two-dimensional physical carrier. Let $r\ge1$, let $\vartheta_1,\ldots,\vartheta_r\in(0,\pi)$ be distinct, and let $w_j>0$ with $\sum_jw_j\le1$. For fixed phases $\varphi_j$, consider
\begin{equation}\label{eq:r69-mixture}
 b_n^{\rm mix}=\tfrac12\left[1+\alpha^n
             \sum_{j=1}^r w_j\cos(n\vartheta_j+\varphi_j)\right].
\end{equation}
Write $m_{\rm mix}(\alpha)$ for its least pure-start stochastic memory.

\begin{theorem}[Spectral mixtures preserve the matched arithmetic exponent]\label{thm:r69-finite-mixture}
For $0<\alpha<1$, the exact memory of \eqref{eq:r69-mixture} satisfies
\begin{equation}\label{eq:r69-mixture-bounds}
 \max\{2r+1,\max_j m(\alpha,\vartheta_j)\}
 \le m_{\rm mix}(\alpha)
 \le 1+\one_{\{\sum_jw_j<1\}}+\sum_{j=1}^r m(\alpha,\vartheta_j).
\end{equation}
All models are realized with the counted prescribed initial label. If every fixed $\vartheta_j/(2\pi)$ has a global badly-approximable certificate, then
$m_{\rm mix}(\alpha)=\Theta((1-\alpha)^{-1/3})$, with constants depending on the finite phase list. The same conclusion holds when some phases are rational and at least one is certified badly approximable. This probability response is physically prepared on dimension at most $2r+1$, and its exact memory is the boundary-optimal observer memory of the same squared task after the standard two-dimensional herald-flag embedding.
\end{theorem}
\begin{proof}
The $2r+1$ modes $1,\alpha e^{\pm i\vartheta_j}$ are pairwise distinct and have nonzero coefficients. Their Hankel rank is $2r+1$ by Vandermonde independence. Every realizing characteristic polynomial vanishes at every nonconstant mode, giving both lower bounds through Theorem~\ref{thm:r69-spectral}.

For each component choose a minimal stochastic realization $(Q_j,d_j)$ starting in its first label. Keep their state sets disjoint and add one initial label $o$ with decoder $b_0^{\rm mix}$. On its first failure, its row enters block $j$ with subprobability row $w_je_1^{\mathsf T}Q_j$. If $w_0=1-\sum_jw_j>0$, its remaining mass enters one absorbing label with decoder $1/2$. Within each block use $Q_j$. Thus after $n\ge1$ failures the output probability is
$w_0/2+\sum_jw_je_1^{\mathsf T}Q_j^nd_j=b_n^{\rm mix}$, while the initial decoder handles $n=0$. No physical call is added and no persistent choice is stored outside the counted labels. This proves the upper. Theorem~\ref{thm:r69-cubic} supplies matched exponents for the finite certified phase list. For a rational phase of order $q$, the matrix $\alpha C_q+(1-\alpha)J$ gives a uniform $q$-label upper, so adding those components does not change the exponent.

For the physical realization, take a direct sum of $r$ qubits, prepare the phase-rotated pure state in block $j$ with weight $w_j$, and use the damped rotation channel in that block. Drop off-diagonal block coherences. If $w_0>0$, add one scalar block of mass $w_0$ with effect $1/2$ and identity evolution. Each diagonal branch is completely positive and their direct sum is trace preserving. The terminal effect is the direct sum of the qubit $Z$ effects and the scalar half effect. Its probability is exactly \eqref{eq:r69-mixture}. Theorem~\ref{thm:r67-embedding} therefore gives the same fixed squared loss and cost $1/p$, and forces this response at the optimum. The observer's independently randomized block selection does not reveal or copy the protected block.
\end{proof}

The three memory statements have different quantifiers: the exact worst-phase law permits the phase to vary with the noise, the cubic-root law fixes any phase with the displayed arithmetic certificate, and the last corollary fixes the older algebraic rotation without such a certificate. Noise-dependent classical memory bounds for quantum processes were already established by Fanizza, Lumbreras and Winter \cite{Fanizza24}; their worst-phase process bound is not a fixed-phase matched law for this response task. The stochastic regions and their Type-0 rigidity are also classical. The present proofs make the task-level reduction and the new arithmetic comparison explicit, rather than claiming earliest priority from a literature search. No statement in this section solves the different squared-sign prediction frontier or the phase-type timing frontier.
